\begin{frame}{$k$ 트레이드오프: 편향과 분산}
  \begin{columns}[T]
    \begin{column}{0.48\textwidth}
      \kb{$k$ 너무 작음} (예: $k=1$)
      \begin{itemize}\small
        \item 노이즈 한 점에도 민감하게 반응
        \item 훈련 데이터를 살짝만 바꿔도 예측이 크게 요동
        \item \textbf{분산(variance)이 크고} $\Rightarrow$
          \kb{과적합}
      \end{itemize}
    \end{column}
    \begin{column}{0.48\textwidth}
      \kb{$k$ 너무 큼} (예: $k=m$, 전체)
      \begin{itemize}\small
        \item 항상 전체 다수결만 $\to$ 지역적 패턴 완전 무시
        \item 분산은 작지만 구조적으로 패턴을 못 잡음
        \item \textbf{편향(bias)이 크고} $\Rightarrow$
          \kb{과소적합}
      \end{itemize}
    \end{column}
  \end{columns}
  \vskip0.8em
  이 관계는 kNN만의 얘기가 아니라 \textbf{지도학습 전체}에 적용되는
  원리. 회귀의 기대 오차는 세 항으로 분해된다:
\end{frame}
